\uuid{auU3}
\titre{Matrice d'une application linéaire/Matrice de passage/Polynômes}
\auteur{hanani ; mzari}
\organisation{exo7}
\licence{CC BY-NC-SA 4.0}
\datecreate{2020-01-20}
\niveau{L1}
\module{Algèbre}
\chapitre{Matrice}
\sousChapitre{Matrice et application linéaire}
\difficulte{2}
\theme{Matrice d'une application linéaire, Matrice de passage, Polynômes}
\exo7qcmid{lille-2.matrice-application.009}

\contenu{
\begin{question}
\qclassification{108.03}
\qtags{motcle=Matrice d'une application linéaire/Matrice de passage/Polynômes}
Dans $\Rr_2[X]$, l'ensemble des polynômes à coefficients réels de degré $\le 2$, on considère la base canonique ${\cal {B}} =\{1,X,X^2\}$ et  la base ${\cal {B'}} =\{P_1,P_2,P_3\}$, où
$P_1=X, P_2=1-X$ et $P_3=(1-X)^2$. On notera $P$ la matrice de passage de la base ${\cal {B}}$ à la base ${\cal {B'}}$ et $Q$ la matrice de passage de la base ${\cal {B'}}$ à la base ${\cal {B}}$.
\vskip1mm
Définition : Soit $E$ un espace vectoriel muni de deux bases ${\cal {B}}$ et ${\cal {B}}'$. La matrice de passage de la base ${\cal {B}}$ à la base  ${\cal {B}}'$ est la matrice de l'identité de $E$ de la base ${\cal {B}}'$ à la base  ${\cal {B}}$. Autrement dit, c'est la matrice dont la jième colonne est constituée des coordonnées du jième vecteur de la base ${\cal {B}'}$ dans la base  ${\cal {B}}$.
\vskip1mm
Quelles sont les assertions vraies ?
\begin{answers}
\bad{$P = \left(\begin{array}{rcc}
1&1&1\\
1&0&-1\\
0&0&1\\
\end{array}\right).$}
\good{$Q = \left(\begin{array}{rcc}
1&1&1\\
1&0&-1\\
0&0&1\\
\end{array}\right).$}
\bad{$Q = \left(\begin{array}{rcc}
0&1&1\\
1&-1&-2\\
0&0&1\\
\end{array}\right).$}
\good{La matrice de l'application identité de $\Rr_2[X]$ de la base ${\cal {B'}}$ à la base ${\cal {B}}$ est :
$$\left(\begin{array}{rcc}
0&1&1\\
1&-1&-2\\
0&0&1\\
\end{array}\right).$$}
\end{answers}
\vskip2mm
\begin{explanations} Par définition, $P$ est la matrice de l'application identité de $\Rr_2[X]$ de la base ${\cal {B'}}$ à la base ${\cal {B}}$ :
$$P= \left(\begin{array}{rcc}
0&1&1\\1&-1&-2\\ 0&0&1\end{array}\right)\quad \mbox{et}\quad
Q = \left(\begin{array}{rcc}1&1&1\\
1&0&-1\\ 0&0&1\end{array}\right).$$
\end{explanations}
\end{question}
}
